\chapter{User Documentation}

\section{Short Problem Statement}
\textit{A Task for Two} is a cooperative two-player puzzle application in which each player has incomplete local information. Neither player can solve the puzzle independently because each side controls a die that contributes to the other player's missing equation value.

The practical user problem solved by the game is coordinated remote puzzle solving, where both players must create/join the same online session, communicate observed values, perform synchronized actions, and complete a shared objective. The level can be finished only when both sides maintain the correct puzzle state long enough for the barrier to open.

\section{Brief Method Overview}
The implemented method combines asymmetric puzzle information with synchronized multiplayer state updates. Each side receives one equation with a hidden value, while each player can rotate exactly one die. Puzzle progress therefore depends on exchanging information and setting values for the counterpart rather than self-solving.

Session connectivity is managed through Unity Relay and host/client Netcode startup. During runtime, each player controls only their own player object. Puzzle state, die lock status, and end-state triggers are synchronized between both players. The solve algorithm requires continuous correctness for a hold duration, which prevents accidental success from random rapid interactions.

The user-facing flow is intentionally linear and explicit: session setup in the main menu, puzzle collaboration in gameplay, then one of two terminal states (victory end screen or connection-lost recovery).

\section{System Requirements}
\subsection*{User Requirements}
The runtime application targets Windows standalone execution with keyboard and mouse input.

\begin{table}[H]
\centering
\caption{Runtime Requirements}
\label{tab:runtime-requirements}
\begin{tabular}{|>{\raggedright\arraybackslash}p{0.2\linewidth}|>{\raggedright\arraybackslash}p{0.7\linewidth}|}
\hline
\textbf{Requirement} & \textbf{Details} \\
\hline
Operating system & Windows (standalone Unity player target). \\
\hline
Input devices & Keyboard and mouse are required for the default control scheme. \\
\hline
Display & Any desktop resolution supported by the monitor and Unity player. \\
\hline
Network & Internet access is required for host/join multiplayer sessions via Relay. \\
\hline
Audio output & Recommended for menu feedback, die roll, barrier, and end-state sound cues. \\
\hline
\end{tabular}
\end{table}

\subsection*{Network Requirements}
\begin{itemize}
\item Both players must run compatible builds of the same project version.
\item The host shares a six-character join code generated at session creation.
\item A stable internet connection is required; disconnects route to a dedicated recovery screen.
\end{itemize}

\subsection*{Development Requirements (optional)}
For development, testing, and running from source, the following tools were used:
\begin{itemize}
\item Unity Hub,
\item Unity Editor \texttt{6000.3.10f1},
\item Git,
\item Git LFS,
\item Visual Studio Code \texttt{1.111.0}.
\end{itemize}

\section{Installation and Startup}
\begin{enumerate}
\item Obtain the project source by cloning the repository:
\begin{verbatim}
git clone https://github.com/hassan-afifi/A_Task_For_Two.git
\end{verbatim}
or
\begin{verbatim}
git clone git@github.com:hassan-afifi/A_Task_For_Two.git
\end{verbatim}

\item Enter the project directory:
\begin{verbatim}
cd A_Task_For_Two
\end{verbatim}

\item Pull large assets managed by Git LFS:
\begin{verbatim}
git lfs install
git lfs pull
\end{verbatim}

\item Open Unity Hub, choose \textbf{Add project from disk}, select the repository folder, and open it with Unity Editor \texttt{6000.3.10f1}.

\item Open the \texttt{MainMenu} scene, press Play, and verify that menu panels, name input, and options UI load correctly.
\end{enumerate}

\begin{figure}[H]
\centering
\includegraphics[width=0.95\textwidth]{images/UnityProject.png}
\caption{Project opened in Unity Editor.}
\label{fig:unity-project-open}
\end{figure}

\section{Controls and Interaction}
The control scheme follows standard first-person keyboard and mouse bindings. Movement and camera look are continuous actions, while interaction and menu actions are event-driven.

\begin{table}[H]
\centering
\caption{Default Keyboard and Mouse Controls}
\label{tab:default-controls}
\begin{tabular}{|>{\raggedright\arraybackslash}p{0.15\linewidth}|>{\raggedright\arraybackslash}p{0.3\linewidth}|>{\raggedright\arraybackslash}p{0.45\linewidth}|}
\hline
\textbf{Action} & \textbf{Default Input} & \textbf{Behavior} \\
\hline
Move & \texttt{W/A/S/D} (or arrow keys) & Moves the local player in first-person space. \\
\hline
Look & Mouse movement & Rotates player body (yaw) and camera pitch. \\
\hline
Jump & \texttt{Space} & Performs jump when grounded and not blocked by crouch/cooldown. \\
\hline
Sprint & \texttt{Left Shift} (hold) & Multiplies movement speed while not crouching. \\
\hline
Crouch & \texttt{Left Ctrl} & Toggles crouch state and updates capsule/camera height. \\
\hline
Interact & \texttt{E} or \texttt{Enter} & Triggers raycast interaction with the assigned puzzle die. \\
\hline
Pause & \texttt{Escape} & Opens or closes the pause menu during gameplay. \\
\hline
Menu click & Mouse left button & Activates UI controls (buttons, sliders, dropdowns). \\
\hline
\end{tabular}
\end{table}

\begin{figure}[H]
\centering
\includegraphics[width=0.95\textwidth]{images/GameView.png}
\caption{First-person gameplay view during active player control.}
\label{fig:first-person-view}
\end{figure}

\section{Gameplay Flow}
\begin{enumerate}
\item Open the main menu and create or join a session. The host uses \textit{Create Game} while the second player enters the shared join code and uses \textit{Join Game}.

\begin{figure}[H]
\centering
\includegraphics[width=0.95\textwidth]{images/MainMenu.png}
\caption{Main menu interface with session creation and join options.}
\label{fig:main-menu-flow}
\end{figure}

\item Move to the puzzle area once both players are connected. Ownership-based spawn and movement control ensures each player controls only their own character.

\begin{figure}[H]
\centering
\includegraphics[width=0.95\textwidth]{images/SecondPlayer.png}
\caption{Both players connected in the gameplay scene.}
\label{fig:both-players-scene}
\end{figure}

\item Read local puzzle symbols and communicate observed values to the other player. Each side sees only partial equation information, so communication is mandatory.

\begin{figure}[H]
\centering
\includegraphics[width=0.95\textwidth]{images/PuzzleA.png}
\caption{Player A view showing partial puzzle information.}
\label{fig:puzzle-a-view}
\end{figure}

\item Rotate your assigned die to provide the value required by your partner's puzzle. The die can be rolled only when interaction raycast conditions are satisfied.

\begin{figure}[H]
\centering
\includegraphics[width=0.95\textwidth]{images/PuzzleB.png}
\caption{Player B view showing partial puzzle information.}
\label{fig:puzzle-b-view}
\end{figure}

\item Keep both dice on the correct values until the hold condition is met. If a wrong value appears or a die is still rolling, hold progress resets.

\begin{figure}[H]
\centering
\includegraphics[width=0.95\textwidth]{images/BarsUp1.png}
\caption{Puzzle solved state with opened barrier.}
\label{fig:puzzle-solved-bars}
\end{figure}

\newpage

\item Move through the opened path to trigger the shared victory sequence and end screen.

\begin{figure}[H]
\centering
\includegraphics[width=0.95\textwidth]{images/EndScreen.png}
\caption{End screen displayed after level completion.}
\label{fig:end-screen-flow}
\end{figure}

\item If a network disconnect occurs before victory, the connection-lost screen is shown instead of normal gameplay continuation.

\begin{figure}[H]
\centering
\includegraphics[width=0.95\textwidth]{images/HostConnectionLost.png}
\caption{Connection-lost screen displayed when a player disconnects during an active session.}
\label{fig:connection-lost-flow}
\end{figure}
\end{enumerate}

\section{Session Validation}
The menu layer applies deterministic input validation before allowing session operations:
\begin{itemize}
\item Player name must contain at least one letter.
\item Join code must contain exactly six characters after trimming.
\item Join code is normalized to uppercase before join request.
\item Create/Join buttons are interactable only when corresponding validation rules pass.
\end{itemize}

These checks reduce invalid network requests and provide immediate feedback before Relay operations start.

\section{Options and HUD}
The options screen applies changes at runtime and stores them in local persistent storage.

\subsection*{Display and Camera Settings}
\begin{table}[H]
\centering
\caption{Display and Camera Options}
\label{tab:display-camera-options}
\begin{tabular}{|>{\raggedright\arraybackslash}p{0.2\linewidth}|>{\raggedright\arraybackslash}p{0.25\linewidth}|>{\raggedright\arraybackslash}p{0.45\linewidth}|}
\hline
\textbf{Setting} & \textbf{Range / Modes} & \textbf{Runtime Behavior} \\
\hline
Display mode & Windowed / Borderless / Fullscreen & Applies fullscreen mode immediately; borderless uses highest-ranked available resolution. \\
\hline
Resolution & Available monitor resolutions & In fullscreen, list is filtered by aspect ratio; in borderless, manual selection is disabled. \\
\hline
Quality level & Unity quality presets & Applies selected quality index through \texttt{QualitySettings}. \\
\hline
Field of view & 60--100 (default 80) & Applies to owner camera and persists. \\
\hline
Sensitivity & 1--100\% (step 1) & Mapped to runtime sensitivity 0.01--1.00 and applied to owner movement. \\
\hline
\end{tabular}
\end{table}

\subsection*{Audio Settings}
\begin{table}[H]
\centering
\caption{Audio Options}
\label{tab:audio-options}
\begin{tabular}{|>{\raggedright\arraybackslash}p{0.25\linewidth}|>{\raggedright\arraybackslash}p{0.1\linewidth}|>{\raggedright\arraybackslash}p{0.55\linewidth}|}
\hline
\textbf{Setting} & \textbf{Range} & \textbf{Runtime Behavior} \\
\hline
Master volume & 0--100\% & Controls master mixer parameter (fallback to listener when needed). \\
\hline
Game SFX volume & 0--100\% & Controls gameplay sound effects channel. \\
\hline
Menu SFX volume & 0--100\% & Controls UI hover/click sound channel. \\
\hline
Game music volume & 0--100\% & Controls gameplay music channel; can be muted by end/disconnect flow. \\
\hline
Menu music volume & 0--100\% & Controls menu music channel. \\
\hline
\end{tabular}
\end{table}

\subsection*{Crosshair and HUD Settings}
\begin{table}[H]
\centering
\caption{Crosshair and HUD Options}
\label{tab:hud-options}
\begin{tabular}{|>{\raggedright\arraybackslash}p{0.25\linewidth}|>{\raggedright\arraybackslash}p{0.2\linewidth}|>{\raggedright\arraybackslash}p{0.45\linewidth}|}
\hline
\textbf{Setting} & \textbf{Range / States} & \textbf{Runtime Behavior} \\
\hline
Crosshair size & 1--10 & Mapped to 5--50 pixel UI size. \\
\hline
Crosshair color & 10 preset colors & Applies immediately to crosshair graphic. \\
\hline
Show FPS & On / Off & Toggles FPS widget visibility. \\
\hline
Show Ping & On / Off & Toggles ping widget visibility. \\
\hline
Show System Clock & On / Off & Toggles clock widget visibility. \\
\hline
\end{tabular}
\end{table}

All options are saved automatically and reloaded on subsequent sessions.

\section{Use Case Diagram}
\begin{figure}[H]
\centering
\includegraphics[width=\textwidth,height=0.92\textheight,keepaspectratio]{images/UseCaseDiagram.png}
\caption{A Task For Two use case diagram.}
\label{fig:use-case-diagram}
\end{figure}

\subsection*{Use Case Diagram Description}
Figure~\ref{fig:use-case-diagram} summarizes user-facing behavior from session initialization to terminal states. The primary actor is \textit{Player}, specialized as \textit{Host Player} and \textit{Client Player}. \textit{Unity Relay Service} appears as an external actor because connectivity depends on Relay allocation and join flows.

The use cases represent the complete operational path:
\begin{itemize}
\item setup actions (set player name, choose character, configure settings),
\item session actions (create session, join session, connect via Relay),
\item runtime actions (control player, pause/resume, interact with puzzle, solve puzzle),
\item terminal actions (complete level, handle connection loss, view end state).
\end{itemize}

\subsection*{Relationship Explanation}
The relationships describe dependency and specialization semantics:
\begin{itemize}
\item Host and client specialize common player capabilities but execute different session-entry actions.
\item Session creation/joining depends on external Relay service interaction.
\item Puzzle interaction includes solve behavior; solve behavior includes completion progression.
\item Completion and disconnect both lead to terminal UI flows.
\end{itemize}

Overall, the model captures the complete user interaction lifecycle, from session initialization to termination states.

\newpage

\section{User Stories}
The following user stories document expected behavior in Given--When--Then format and are aligned with implemented runtime flow.

\subsection*{Main Menu and Session Flow}
\begin{longtable}{|>{\raggedright\arraybackslash}p{0.18\linewidth}|>{\raggedright\arraybackslash}p{0.22\linewidth}|>{\raggedright\arraybackslash}p{0.15\linewidth}|>{\raggedright\arraybackslash}p{0.35\linewidth}|}
\caption{User Stories: Main Menu and Session Flow}\label{tab:userstories-main-menu}\\
\hline
\textbf{I want to} & \textbf{Given} & \textbf{When} & \textbf{Then} \\
\hline
Open the options menu & The main menu is visible & I click \textit{Options} & The application displays the options screen \\
\hline
Close the options menu & The options screen is visible & I click \textit{Back} or \textit{Close} & The interface returns to the main menu panel \\
\hline
Create a host session & I entered a valid player name & I click \textit{Create Game} & Hosting starts and gameplay loading begins \\
\hline
Join an existing session & I entered a valid player name and join code & I click \textit{Join Game} & Client connection starts and the player joins the session \\
\hline
Avoid invalid session actions & Name or code input is invalid & I stay on the menu screen & Create and join actions remain unavailable \\
\hline
Exit the application & I am in a menu screen & I click \textit{Exit} & The application closes \\
\hline
\end{longtable}

\newpage

\subsection*{Puzzle Interaction}
\begin{longtable}{|>{\raggedright\arraybackslash}p{0.20\linewidth}|>{\raggedright\arraybackslash}p{0.20\linewidth}|>{\raggedright\arraybackslash}p{0.20\linewidth}|>{\raggedright\arraybackslash}p{0.30\linewidth}|}
\caption{User Stories: Puzzle Interaction}\label{tab:userstories-puzzle}\\
\hline
\textbf{I want to} & \textbf{Given} & \textbf{When} & \textbf{Then} \\
\hline
See my local puzzle information & The gameplay scene is active & I approach the puzzle area & The interface shows my local symbols and target equation values \\
\hline
Rotate my assigned die & I am close to the die and interaction is available & I press the interact key & The die performs one roll step \\
\hline
Provide the other player's missing value & Both players are communicating puzzle data & I set my die to the requested face & The shared puzzle state advances toward a valid solution \\
\hline
Solve the puzzle cooperatively & Both dice are on correct values & The values remain stable for the hold duration & The puzzle is marked as solved \\
\hline
Open the bars after a valid solve & The puzzle was solved & The solve state is confirmed & The barrier opening animation starts \\
\hline
Prevent post-solve random changes & The puzzle is solved & I try to rotate a die again & Further die-roll attempts are ignored while dice stay locked \\
\hline
\end{longtable}

\newpage

\subsection*{Pause and Options During Gameplay}
\begin{longtable}{|>{\raggedright\arraybackslash}p{0.16\linewidth}|>{\raggedright\arraybackslash}p{0.22\linewidth}|>{\raggedright\arraybackslash}p{0.16\linewidth}|>{\raggedright\arraybackslash}p{0.36\linewidth}|}
\caption{User Stories: Pause and Options During Gameplay}\label{tab:userstories-pause-options}\\
\hline
\textbf{I want to} & \textbf{Given} & \textbf{When} & \textbf{Then} \\
\hline
Open the pause menu & Gameplay is active and no blocking screen is shown & I press the pause input & The pause menu overlay appears \\
\hline
Resume gameplay & The pause menu is open & I click \textit{Continue} or close the menu & The overlay closes and control is restored \\
\hline
Open options from pause & The pause menu is open & I click \textit{Options} & The options screen opens \\
\hline
Apply changed options & The options screen is open & I change a setting value & The new setting is applied and stored persistently \\
\hline
Keep gameplay blocked while menu is open & Pause or options UI is open & I attempt world interaction & World interaction stays blocked until the UI is closed \\
\hline
\end{longtable}

\newpage

\subsection*{End and Recovery Flow}
\begin{longtable}{|>{\raggedright\arraybackslash}p{0.20\linewidth}|>{\raggedright\arraybackslash}p{0.25\linewidth}|>{\raggedright\arraybackslash}p{0.15\linewidth}|>{\raggedright\arraybackslash}p{0.30\linewidth}|}
\caption{User Stories: End and Recovery Flow}\label{tab:userstories-end-recovery}\\
\hline
\textbf{I want to} & \textbf{Given} & \textbf{When} & \textbf{Then} \\
\hline
Finish the level for both players & The bars are open and at least one player reaches the end trigger & The trigger is entered & The end screen is shown for both players \\
\hline
Stop gameplay actions on victory & The end screen is shown & I try to move or interact & Gameplay controls remain disabled \\
\hline
Return to main menu after victory & The end screen is shown & I click \textit{Main Menu} & The main menu scene loads \\
\hline
Handle disconnect safely & A client disconnect occurs before victory & Network callback is received & The connection-lost screen appears \\
\hline
Exit from the connection-lost screen & The connection-lost screen is open & I click \textit{Exit} & The application closes \\
\hline
\end{longtable}
